% ---- pg-548 （书页 536） ----

\noindent\textbf{例 12.6}\quad (11 种语言的完全连接聚类)

在例 12.4 中我们给出了 11 种语言的距离阵. 将完全连接聚类算法应用于这个距离阵, 得出的树图见\hyperref[fig:12.7]{图 12.7}.

比较\hyperref[fig:12.7]{图 12.7} 和\hyperref[fig:12.4]{图 12.4}, 我们看到两种分层聚类法都产生出英语—挪威语—丹麦语组和法语—意大利语—西班牙语组. 在中间水平上, 波兰语被合并到法语—意大利语—西班牙语组. 此外, 两种方法都在倒数第二步才将匈牙利语和芬兰语合成一组.

\begin{mvfigure}
\begin{tikzpicture}[x=0.92cm,y=0.42cm,line width=0.5pt]

% ---------- 叶结点的竖线 ----------
\draw (0,0) -- (0,1.91);
\draw (1,0) -- (1,0.86);
\draw (2,0) -- (2,0.86);
\draw (3,0) -- (3,5.93);
\draw (4,0) -- (4,0.86);
\draw (5,0) -- (5,0.86);
\draw (6,0) -- (6,1.91);
\draw (7,0) -- (7,4.91);
\draw (8,0) -- (8,6.86);
\draw (9,0) -- (9,8.07);
\draw (10,0) -- (10,8.07);

% ---------- 合并的横线 ----------
\draw (1,0.86) -- (2,0.86);
\draw (4,0.86) -- (5,0.86);
\draw (0,1.91) -- (1.5,1.91);
\draw (4.5,1.91) -- (6,1.91);
\draw (5.25,4.91) -- (7,4.91);
\draw (0.75,5.93) -- (3,5.93);
\draw (1.875,6.86) -- (8,6.86);
\draw (9,8.07) -- (10,8.07);
\draw (4,9.02) -- (9.5,9.02);

% ---------- 各聚类向上的竖线 ----------
\draw (1.5,0.86) -- (1.5,1.91);
\draw (0.75,1.91) -- (0.75,5.93);
\draw (4.5,0.86) -- (4.5,1.91);
\draw (5.25,1.91) -- (5.25,4.91);
\draw (6.125,4.91) -- (6.125,6.86);
\draw (1.875,5.93) -- (1.875,6.86);
\draw (4,6.86) -- (4,9.02);
\draw (9.5,8.07) -- (9.5,9.02);

% ---------- 坐标轴 ----------
\draw (-0.85,0) -- (-0.85,10.3);
\foreach \t in {0,2,4,6,8,10}{%
  \draw (-0.85,\t) -- (-0.62,\t);%
  \node[left,font=\small] at (-0.95,\t) {\t};%
}
\node[rotate=90,font=\small] at (-1.75,5) {距离};
\node[font=\small] at (5,-2.1) {语言};

% ---------- 叶结点的圆点与标记 ----------
\foreach \i/\l in {0/E,1/N,2/Da,3/G,4/Fr,5/I,6/Sp,7/P,8/Du,9/H,10/Fi}{%
  \fill (\i,0) circle (0.055);%
  \node[below,font=\small,inner sep=1.5pt] at (\i,0) {\l};%
}

\end{tikzpicture}
\phantomsection\label{fig:12.7}\mvcap{图 12.7\quad 11 种语言中前 10 个数词之间距离的完全连接树图}
\end{mvfigure}

然而, 两种方法处理德语和荷兰语的结果不同. 单连接法在中等距离上将德语与荷兰语合并, 而且它们在最后的大合并之前始终是一个单独的聚类. 完全连接法在中间水平上将德语并入英语—挪威语—丹麦语组. 荷兰语保持为单独一组, 直到在较高的距离水平上才与英语—挪威语—丹麦语—德语组及法语—意大利语—西班牙语—波兰语组合并成一组. 完全连接的最终合并涉及两个聚类, 而单连接的最终合并涉及三个聚类. \hfill$\blacksquare$

\medskip
\noindent\textbf{例 12.7}\quad (用完全连接法对变量进行聚类)

从 22 家美国公用事业公司搜集到的 1975 年的数据列于\hyperref[tab:12.4]{表 12.4}. 虽然将公司进行聚类更令人感兴趣, 不过在这里我们要看看完全连接法是如何用于对变量进行聚类的.

\begin{mvtable}
\phantomsection\label{tab:12.4}\mvtabcap{表 12.4\quad 公用事业公司数据(1975 年)}
\footnotesize
\begin{tabular}{lrrrrrrrr}
\toprule
 & \multicolumn{8}{c}{变\quad 量}\\
\cmidrule(lr){2-9}
公\quad 司 & \multicolumn{1}{c}{$X_1$} & \multicolumn{1}{c}{$X_2$} & \multicolumn{1}{c}{$X_3$} & \multicolumn{1}{c}{$X_4$} & \multicolumn{1}{c}{$X_5$} & \multicolumn{1}{c}{$X_6$} & \multicolumn{1}{c}{$X_7$} & \multicolumn{1}{c}{$X_8$}\\
\midrule
1. 亚利桑那公共服务公司 & 1.06 & 9.2 & 151.0 & 54.4 & 1.6 & 9 077.0 & 0 & 0.628\\
2. 波士顿爱迪生公司 & 0.89 & 10.3 & 202.0 & 57.9 & 2.2 & 5 088.0 & 25.3 & 1.555\\
3. 中央路易斯安那电气公司 & 1.43 & 15.4 & 113.0 & 53.0 & 3.4 & 9 212.0 & 0 & 1.058\\
4. 联邦爱迪生公司 & 1.02 & 11.2 & 168.0 & 56.0 & 0.3 & 6 423.0 & 34.3 & 0.700\\
5. 联合爱迪生公司(纽约) & 1.49 & 8.8 & 192.0 & 51.2 & 1.0 & 3 300.0 & 15.6 & 2.044\\
6. 佛罗里达电力与电灯公司 & 1.32 & 13.5 & 111.0 & 60.0 & $-2.2$ & 11 127.0 & 22.5 & 1.241\\
7. 夏威夷电气公司 & 1.22 & 12.2 & 175.0 & 67.6 & 2.2 & 7 642.0 & 0 & 1.652\\
8. 爱达荷电力公司 & 1.10 & 9.2 & 245.0 & 57.0 & 3.3 & 13 082.0 & 0 & 0.309\\
9. 肯塔基公用事业公司 & 1.34 & 13.0 & 168.0 & 60.4 & 7.2 & 8 406.0 & 0 & 0.862\\
10. 麦迪逊天然气与电力公司 & 1.12 & 12.4 & 197.0 & 53.0 & 2.7 & 6 455.0 & 39.2 & 0.623\\
11. 内华达电力公司 & 0.75 & 7.5 & 173.0 & 51.5 & 6.5 & 17 441.0 & 0 & 0.768\\
12. 新英格兰电气公司 & 1.13 & 10.9 & 178.0 & 62.0 & 3.7 & 6 154.0 & 0 & 1.897\\
13. 北方电力公司 & 1.15 & 12.7 & 199.0 & 53.7 & 6.4 & 7 179.0 & 50.2 & 0.527\\
14. 俄克拉何马天然气与电力公司 & 1.09 & 12.0 & 96.0 & 49.8 & 1.4 & 9 673.0 & 0 & 0.588\\
15. 太平洋天然气与电力公司 & 0.96 & 7.6 & 164.0 & 62.2 & $-0.1$ & 6 468.0 & 0.9 & 1.400\\
16. 普盖桑德电力与电灯公司 & 1.16 & 9.9 & 252.0 & 56.0 & 9.2 & 15 991.0 & 0 & 0.620\\
\bottomrule
\end{tabular}
\end{mvtable}
